\documentclass[11pt,a4paper]{article}
\usepackage[margin=25mm]{geometry}
\usepackage{amsmath,amssymb,amsthm,microtype}
\usepackage[hidelinks]{hyperref}
\newtheorem{theorem}{Theorem}\newtheorem{corollary}[theorem]{Corollary}
\newcommand{\E}{\mathbb E}\newcommand{\Pp}{\mathbb P}
\title{A square-root tangent removes the compound-Poisson branch}
\author{}\date{29 October 2026}
\begin{document}\maketitle

The five-count argument already controls the full signed discrepancy at
moderate noise. Its lower endpoint can be reduced from $.075$ to $.03$
without a new probability estimate. The only change is to use a better
affine upper bound for $\sqrt\lambda$ in the curvature error. Four rational
endpoint inequalities then replace the 291-cell negative-phase certificate
in the selected $n\ge200$ right-strip proof. The older $n\ge42$ theorem
retains its larger index domain and is preserved.

\begin{theorem}\label{thm:new}
Let $n\ge200$ be an integer, $p\in[2/5,53/100]$, $q=1-p$ and $v=pq$.
Let $Y=B+U$, where $B$ is Bernoulli of parameter $p$ and $\E(U\mid B)=0$.
Write $U_j$ for the conditional noise given $B=j$, and assume
\[
 -3/5\le U_0\le a_0,\qquad -a_1\le U_1\le3/5,\qquad
 a_0,a_1\ge0,\qquad a_0+a_1\le9/125.
\]
Set $s=\E U^2$, $\lambda=ns$, $X=(Y-p)/\sqrt{v+s}$ and
$\rho=\E|X|^3$. If $3/100\le\lambda\le9/25$, then independent copies
$X_1,\ldots,X_n$ satisfy, for every raw threshold $|H|\le46/25$,
\begin{equation}\label{eq:target}
 \mathcal E:=\sqrt n\{\Pp(X_1+\cdots+X_n\le H)-\Phi(H/\sqrt n)\}
             -C_E\rho<-\lambda/100,
 \qquad C_E=\frac{3+\sqrt{10}}{6\sqrt{2\pi}}.
\end{equation}
\end{theorem}

The conclusion covers both signs of $H$, all phases and all threshold
atoms, but concerns the positive signed discrepancy. Acquiring the strip
for a maximizing witness, and applying it to a reflected law, are separate
steps. This is not an unrestricted cutoff theorem.

\section{The retained probability estimates}
The relevant ancestors are the v13 weighted crossing identity, the v14
summable-curvature lemma, and the v15 five-count reserve. Their precise
source locations are listed below. We give all estimates needed to join them.
Write $t=np+w$, where $w=H\sqrt{v+s}$, and $k=\lfloor t\rfloor$,
$u=t-k\in[0,1)$. Then $|w|<231/250$ and the crossing identity reads
\begin{equation}\label{eq:cross}
 \Pp\Big(\sum X_i\le H\Big)
 =(1-u)F_B(k)+uF_B(k+1)+\mathcal C(t)-\sum_{i=1}^n\pi_i g_i(t),
\end{equation}
where $\pi_i$ and $F_B$ are binomial masses and CDF values, and $g_i\ge0$.
The crossing is interior since
$n(p-a_0)>|w|$ and $n(q-a_1)>|w|$.

For the five indices $k-1,\ldots,k+3$, one has $|i-np|<4$.
Finite Stirling on the segment $[.38,.55]$, whose variance product exceeds
$.23$, gives
\[
 \sqrt n\pi_i>
 \frac{199}{250}\left(1-\frac{8+1/12}{200(23/100)}\right)>
 \frac{13}{20}.
\]
The five-count reserve is obtained by putting $(2/5)\sin^2(\pi x)$,
restricted to $[0,1]$, below the crossing kernel. Four centered replacements
freeze the conditional sum at the middle count. Their total cost is at most
$10\pi^2\lambda/n$ before the factor $2/5$. The complete omitted-translate
tail is bounded by the fourth moment divided by $16$.
With $\alpha=39/40$, $b=203/200$, $z=3\alpha$ and $x=z\lambda$, this gives
\begin{equation}\label{eq:W}
 \sqrt n\sum_i\pi_i g_i\ge W(\lambda):=
 \frac{13}{50}\left\{
 \frac{x-x^2/2+x^3/6-x^4/24}{2}
 -\frac{(9/25)b\lambda+3b^2\lambda^2}{16}
 -\frac{10(22/7)^2}{200}\lambda\right\}.
\end{equation}
Here the frozen variance lies between $(1-5/n)\lambda$ and
$(1+3/n)\lambda$. Independence and the support bound give its fourth
moment at most $(9/25)\tau+3\tau^2$, and its characteristic function at
$2\pi$ has modulus at most $e^{-3\tau}$. Finally
$1-e^{-x}\ge x-x^2/2+x^3/6-x^4/24$ for every $x\ge0$.
None of these estimates needs $\lambda\ge.075$.

For the summable curvature put
\[
 c=(np-q-t)_+,\qquad T=\frac{100}{47}\lambda,\qquad
 M=\frac9{25}T+3T^2.
\]
On the increasing side of the binomial mass the actual conditional variance
is at most $\lambda/q\le T$, and its fourth moment is at most $M$.
The retained curvature lemma, with mode coefficient $41/100$, yields
\begin{equation}\label{eq:curvature}
 \sqrt n\mathcal C(t)\le
 \frac{41/100}{nv^{3/2}}
 \left\{\frac T4+c\sqrt T+
 \frac{27M}{256}\left(\frac53+\frac{5c}{4}\right)\right\}.
\end{equation}
Here $c<23/50$ if $H\le0$, and $c=0$ if $H\ge0$.
The bound on the binomial mode and the hypotheses $c\le1$ and $v^{3/2}>117/1000$
are unchanged from the v15 derivation. All count tails in this lemma are
summed analytically; no finite truncation is being used.

\section{The tangent and the signed normal errors}
The useful replacement is simply
\begin{equation}\label{eq:tangent}
 \sqrt\lambda\le\frac53\lambda+\frac3{20},\qquad
 \sqrt T<\frac32\sqrt\lambda\le\frac52\lambda+\frac9{40}.
\end{equation}
The first inequality is equivalent to
$(\sqrt\lambda-3/10)^2\ge0$; the second uses $100/47<9/4$.
Define the quadratic polynomial
\[
 S_*(\lambda,c)=\frac T4+c\left(\frac52\lambda+\frac9{40}\right)
        +\frac{27M}{256}\left(\frac53+\frac{5c}{4}\right).
\]
All its coefficients are nonnegative for the two choices of $c$ used below.

Schulz's strict binomial theorem \cite[Theorem 1]{Schulz} bounds both integer
anchors in \eqref{eq:cross} by $C_E\rho_B$, where
$\rho_B=(1-2v)/\sqrt v$. Convex interpolation preserves the strict bound.
For $G(y)=\Phi((y-np)/\sqrt{nv})$, the signed interpolation remainder is
at most
\[
 \sqrt n\{(1-u)G(k)+uG(k+1)-G(t)\}
 \le\frac{\phi(0)}{2nv^{3/2}}u(1-u)(u-w)_+.
\]
When $H\le0$ the variance change from $v$ to $v+s$ is favorable.
Using $u(1-u)\le1/4$ and $u^2(1-u)\le4/27$, the total normal cost is
at most
\[
 \frac{2/5}{nv^{3/2}}\left(\frac{231}{2000}+\frac2{27}\right).
\]
When $H\ge0$, its bound is instead
\[
 \frac{2/5}{nv^{3/2}}\left(\frac2{27}+\frac{231}{500}\lambda\right).
\]
The second term comes from $s=\lambda/n$; it has no missing normalization.
The fixed signs on the two conditional supports give
\[
 (v+s)^{3/2}(\rho-\rho_B)
 \ge s\{87/250-(13/16)s\}>0,
\]
since $s\le9/5000$. Thus the actual cubic penalty is favorable.

Combining these facts with \eqref{eq:W}--\eqref{eq:tangent} gives
$\mathcal E<-W(\lambda)+E_\pm(\lambda)/n$, where
\begin{align*}
 E_-(\lambda)&=
 \frac{(41/100)S_*(\lambda,23/50)+(2/5)(231/2000+2/27)}{117/1000},\\
 E_+(\lambda)&=
 \frac{(41/100)S_*(\lambda,0)+(2/5)(2/27+231\lambda/500)}{117/1000}.
\end{align*}
The entire improvement is in $E_-$; the positive branch is unchanged.

\section{Four exact comparisons, with no parameter cells}
Put $P_\pm=W-E_\pm/200-\lambda/100$. Both $E_\pm$ are convex
quadratics with positive coefficients, while
\[
 W''(\lambda)=-\frac{13}{50}
 \left\{\frac{z^2}{4}\bigl[(z\lambda-1)^2+1\bigr]
                         +\frac{3b^2}{8}\right\}<0.
\]
Hence $P_\pm$ are strictly concave on the real line. Exact rational
arithmetic gives the following lower bounds:
\[
\begin{array}{c|cc}
 &\lambda=3/100&\lambda=9/25\\\hline
 P_-&>461/10^6&>885/10^6\\
 P_+&>4654/10^6&>11247/10^6.
\end{array}
\]
Concavity proves positivity on the whole closed interval. Since $E_\pm>0$
and $n\ge200$, this proves Theorem~\ref{thm:new}.
The companion checker reconstructs these polynomials using exact fractions
and checks the displayed second-derivative identity coefficient by coefficient.
No numerical optimizer, parameter mesh or compound-Poisson computation is used.

\begin{corollary}\label{cor:union}
Under the same support, centering, probability and threshold hypotheses,
for $n\ge200$ and $0<\lambda\le9/25$ one has
$\mathcal E<-\lambda/100$. At $\lambda=0$ one has $\mathcal E<0$.
\end{corollary}
\begin{proof}
For $0<\lambda\le3/100$, use the v11 direct comparison at its shifted
integer anchor and $C=C_E>2/5$, followed by the strict binomial theorem.
Its phase interval $[-34/125,91/125]$ has length one, so an anchor always
exists. For $\lambda\ge3/100$, apply Theorem~\ref{thm:new}. The common
endpoint is included in both domains. When $\lambda=0$, both centered
conditional noises vanish and the strict binomial theorem applies directly.
\end{proof}

\section*{Exact dependency changes and limitations}
All paths below begin with \path{mathematics/}.
\begin{sloppypar}
\begin{itemize}
\item Retain \path{simplification-v15/2026-10-28/loss-direct.tex}, the
five-count theorem and its proof; \path{simplification-v13/2026-10-28/crossing.tex},
the weighted identity; and \path{simplification-v14/2026-10-28/geometry-slack.tex},
the summable-curvature lemma. The finite-Stirling and signed normal estimates
are the displayed ones in \path{simplification-v15/2026-10-28/right-direct.tex}.
\item Retain \path{simplification-v11/2026-10-28/counts-v11.tex}, its first
theorem, for positive noise through $.03$; its scalar Poisson-tail argument
remains a dependency. Retain the external binomial theorem.
\item In the $n\ge200$ right-strip union of
\path{simplification-v15/2026-10-28/join-review.md}, replace the $.03$--$.075$
phase split by Theorem~\ref{thm:new}. Neither
\path{simplification-v12/2026-10-28/compound-proof.tex} nor the separate
\path{simplification-v12/2026-10-28/outer.tex} is then needed on that path.
Thus the 291-cell allocation certificate, compound-Poisson evaluator and
its complete-tail and box-coverage subarguments leave this dependency path.
\end{itemize}
\end{sloppypar}
This does not replace the $n\ge42$ anchor-comparison theorem at its original
scope. In particular no statement for $42\le n<200$ is inferred. The
acquired-witness $n\ge700$ union uses only $n\ge200$ on its right side, so
this replacement suffices there. Left-strip estimates, coarse Fourier entry,
support acquisition, reflection, attainment and the genuine finite base are
unchanged. No existing source or certificate is removed.

The old tangent $\sqrt\lambda\le\lambda+1/4$ gives
$P_-(.03)<0$ at $n=200$; that is an obstruction to those sufficient bounds,
not a counterexample to the probability theorem. The present tangent also
cannot be extrapolated to zero noise: its scalar $P_-(0)<0$. This explains
why the v11 small-noise argument remains. Actual-law checks and these
negative controls are recorded separately from continuous-domain acceptance.

\begin{thebibliography}{1}
\bibitem{Schulz}J. Schulz, \emph{The optimal Berry--Esseen constant in the
binomial case}, doctoral dissertation, Universit\"at Trier, 2016.
\url{https://d-nb.info/1197702695/34}.
\end{thebibliography}
\end{document}
